\section{A finite fracture construction}\label{sec:fracture}

This section isolates the topological assembly.  Its inputs are a
height-two equivalence, a compatible height-one and rational pair of
maps, and a basis for the final rational quotient.  The construction
produces the stages themselves and computes every connecting map.
It does not use the coefficient calculations of
\Cref{h3:sec:ordinary}.

\subsection{Module localizations and rotation signs}

Put \(Q=L_0S\simeq H\Qp\) and \(J=K(1)\vee K(2)\).
We work in \(\Mod_S\).  Homological Bousfield localization \(L_E\)
lifts to this category with the same underlying spectrum, whether
or not it is smashing.  Indeed,
\(S^{\wedge n}\wedge M\to S^{\wedge n}\wedge L_EM\) is an
\(E\)-equivalence for every \(n\).  Mapping into the local
spectrum \(L_EM\) extends the action maps and their coherent unit
and associativity homotopies.  The free-module bar resolution
then gives, for every local \(S\)-module \(N\),
\begin{equation}\label{eq:module-localization}
 \Map_{\Mod_S}(L_EM,N)\simeq\Map_{\Mod_S}(M,N).
\end{equation}
Thus the units below are module maps.  The forgetful functor creates
limits and detects equivalences, so the spectral fracture pullbacks
are also pullbacks of \(S\)-modules.

Rationalization is extension of scalars:
\(Q\otimes_SM\simeq L_0M\).  In particular, for a rational
module \(D\),
\begin{equation}\label{eq:rational-adjunction}
 \Map_{\Mod_S}(M,D)
   \simeq\Map_{\Mod_Q}(L_0M,D).
\end{equation}
A map from a finite sum of shifts of \(Q\) is determined up to
homotopy by the images of its module generators.  A homotopy between
the rational maps in \eqref{eq:rational-adjunction} gives the
corresponding \(S\)-module homotopy after precomposition by
\(M\to L_0M\).

The height-two fracture square for an \(E(2)\)-local module \(B\)
is
\begin{equation}\label{eq:height-two-fracture}
 B\simeq L_{K(2)}B
       \mathbin{\times}_{L_1L_{K(2)}B}L_1B.
\end{equation}
For an \(E(1)\)-local module the corresponding square uses \(K(1)\)
and \(L_0\) \cite[Section~2.4, diagram~(2.19)]{BB}.
Consequently an \(E(2)\)-local, \(K(2)\)-acyclic module is
\(E(1)\)-local, and an \(E(2)\)-local module acyclic for both
\(K(1)\) and \(K(2)\) is rational.  There is also a fracture
square separating all positive heights:
\begin{equation}\label{eq:positive-fracture}
 B\simeq L_JB\mathbin{\times}_{L_0L_JB}L_0B.
\end{equation}
To see this, \(L_JB\) is \(E(2)\)-local because
\(\langle E(2)\rangle=\langle K(0)\vee J\rangle\).
The fiber \(N=\fib(B\to L_JB)\) is therefore \(E(2)\)-local
and \(J\)-acyclic, hence rational.  Exactness of \(L_0\) identifies
it with \(\fib(L_0B\to L_0L_JB)\), which proves
\eqref{eq:positive-fracture}.

For \(f:A\to B\) and \(C=\cofib(f)\), we fix the
cofiber-rotation convention
\[
 A\xrightarrow{f}B\longrightarrow C\longrightarrow\Sigma A
 \xrightarrow{-\Sigma f}\Sigma B.
\]
If \(i:N\to A\) is the fiber of \(A\to D\), the induced
identification \(\cofib(A\to D)\simeq\Sigma N\) therefore makes
the boundary to \(\Sigma A\) equal to \(-\Sigma i\).

\begin{lemma}[Boundary of a localization comparison]\label{lem:boundary}
Let \(a:A\to B\) be a map in a stable category, and let \(L\) be
an exact localization for which \(La\) is an equivalence.
Choose an inverse of \(La\) and its inverse homotopies, and put
\[
 \tau:B\longrightarrow LB\xrightarrow{(La)^{-1}}LA,\qquad
 N=\fib(A\longrightarrow LA).
\]
The inverse homotopy makes \(\tau\) a map under \(A\), and
therefore gives
\[
 \bar\tau:\cofib(a)\longrightarrow\cofib(A\to LA)
                         \simeq\Sigma N.
\]
If \(i:N\to A\) is the fiber inclusion, then the connecting map
of the cofiber sequence of \(a\) is
\begin{equation}\label{eq:boundary-lemma}
 \partial_a=(-\Sigma i)\bar\tau.
\end{equation}
If \(M\) is another exact localization and \(N\) is \(M\)-local,
the middle cofiber can equivalently be written
\(\cofib(MA\to MLA)\), using its localization identification.
\end{lemma}

\begin{proof}
The under-\(A\) homotopy supplies a map of cofiber sequences
\[
\begin{tikzcd}[column sep=large]
A \arrow[r,"a"] \arrow[d,equal]
  & B \arrow[r] \arrow[d,"\tau"]
  & \cofib(a) \arrow[r,"\partial_a"] \arrow[d,"\bar\tau"]
  & \Sigma A \arrow[d,equal]\\
A \arrow[r]
  & LA \arrow[r]
  & \Sigma N \arrow[r,"-\Sigma i"]
  & \Sigma A.
\end{tikzcd}
\]
The right square is \eqref{eq:boundary-lemma}; its sign is the
rotation sign fixed above.  If \(N\) is \(M\)-local, so is
\(\cofib(A\to LA)\simeq\Sigma N\).  Its \(M\)-localization
map is then an equivalence, and exactness identifies the localized
cofiber with \(\cofib(MA\to MLA)\).
\end{proof}

\subsection{The stages and their boundaries}

Use the following ordered blocks:
\[
\begin{aligned}
 W&=W_1\vee W_2=L_2S\vee\Sigma^{-1}L_2S,\\
 U_{\mathrm{mid}}&=U_3\vee U_4=\Sigma^{-3}L_1S\vee\Sigma^{-4}L_1S,\\
 V&=V_5\vee V_6\vee V_7\vee V_8\\
  &=\Sigma^{-5}Q\vee\Sigma^{-6}Q
       \vee\Sigma^{-8}Q\vee\Sigma^{-9}Q.
\end{aligned}
\]

\begin{proposition}[Ordered fracture filtration]\label{prop:fracture}
Let \(X\) be an \(E(2)\)-local \(S\)-module.  Suppose the
following module maps and homotopies are given.
\begin{enumerate}[label=\textup{(\roman*)}]
\item A map \(w:W\to X\) for which \(L_{K(2)}w\) is an
  equivalence.  Put \(Y=\cofib(w)\), with quotient \(q:X\to Y\).
\item An equivalence
  \(\kappa:L_{K(1)}U_{\mathrm{mid}}\xrightarrow{\simeq}L_{K(1)}Y\), a rational
  map \(r:L_0U_{\mathrm{mid}}\to L_0Y\), and a specified homotopy of \(Q\)-module
  maps between
  \begin{equation}\label{eq:compatibility}
  \begin{gathered}
  L_0U_{\mathrm{mid}}\longrightarrow L_0L_{K(1)}U_{\mathrm{mid}}
          \xrightarrow{L_0\kappa}L_0L_{K(1)}Y,\\
  L_0U_{\mathrm{mid}}\xrightarrow rL_0Y\longrightarrow L_0L_{K(1)}Y.
  \end{gathered}
  \end{equation}
\item A map \(b:V\to L_0X\) for which
  \begin{equation}\label{eq:quotient-basis}
  t_0:V\xrightarrow bL_0X\xrightarrow{L_0q}L_0Y
                 \longrightarrow\cofib(r)
  \end{equation}
  is an equivalence.
\end{enumerate}
Then there is a filtration in \(E(2)\)-local \(S\)-modules
\[
 0=F_0\longrightarrow F_1\longrightarrow\cdots\longrightarrow F_8
 \xrightarrow[\simeq]{e}X
\]
with ordered cofibers \(W_1,W_2,U_3,U_4,V_5,V_6,V_7,V_8\).
The composite \(F_2=W\to F_8\xrightarrow eX\) is \(w\).
\end{proposition}

\begin{proof}
\emph{Constructing the middle map.}
The cofiber \(Y\) is \(E(2)\)-local and \(K(2)\)-acyclic by
\textup{(i)}, hence \(E(1)\)-local.  Exactness identifies it
naturally with \(\cofib(L_1w)\).  Its height-one fracture square is
\[
 Y\simeq L_{K(1)}Y\mathbin{\times}_{L_0L_{K(1)}Y}L_0Y.
\]
The two proposed maps from \(U_{\mathrm{mid}}\) to its corners are
\[
 U_{\mathrm{mid}}\longrightarrow L_{K(1)}U_{\mathrm{mid}}\xrightarrow{\kappa}L_{K(1)}Y,
 \qquad
 U_{\mathrm{mid}}\longrightarrow L_0U_{\mathrm{mid}}\xrightarrow rL_0Y.
\]
By \eqref{eq:rational-adjunction}, the specified rational homotopy
in \eqref{eq:compatibility} identifies their composites to the
overlap as \(S\)-module maps.  The pullback therefore gives
\begin{equation}\label{eq:middle-map}
 h:U_{\mathrm{mid}}\longrightarrow Y,\qquad L_{K(1)}h=\kappa,\quad L_0h=r.
\end{equation}
Here each equality includes the homotopy obtained from the
corresponding pullback projection: localizing that homotopy and
using idempotence gives the displayed identification.  In
particular, \(h\) is a \(K(1)\)-equivalence.

\emph{The first four stages.}
Let \(\partial_w:Y\to\Sigma W\) be the boundary of \(w\), and put
\[
 d=\partial_wh:U_{\mathrm{mid}}\longrightarrow\Sigma W,\qquad
 d_i=d|_{U_i}\quad(i=3,4).
\]
Define
\begin{equation}\label{eq:first-stages}
 F_0=0,\quad F_1=W_1,\quad F_2=W,\quad
 F_3=X\mathbin{\times}_YU_3,\quad
 F_4=X\mathbin{\times}_YU_{\mathrm{mid}}.
\end{equation}
The first nonzero transition is the wedge inclusion \(W_1\to W\).
The identification \(X\times_Y0\simeq W\) makes the next transitions
the pullbacks of \(0\to U_3\to U_3\vee U_4\).
In a stable category a pullback square is also a pushout square.
Their cofibers are consequently \(U_3\) and \(U_4\), and
\[
 F_3\simeq\fib(d_3:U_3\to\Sigma W),\qquad
 F_4\simeq\fib(d:U_{\mathrm{mid}}\to\Sigma W).
\]
Naturality of the cofiber sequence for each pullback gives
\begin{equation}\label{eq:first-boundaries}
 \partial_1=0,\qquad \partial_2=0,\qquad
 \partial_3=d_3,\qquad
 \partial_4=\Sigma(F_2\to F_3)d_4.
\end{equation}
For the last formula, the quotient of the pullback square for
\(U_3\to U_3\vee U_4\) identifies its quotient with \(U_4\).
There is a map of cofiber sequences from
\(W\to F_4\to U_{\mathrm{mid}}\) to
\(F_3\to F_4\to U_4\), whose components are
\(W\to F_3\), the identity of \(F_4\), and the projection
\(\operatorname{pr}_4:U_{\mathrm{mid}}\to U_4\).
Its boundary square gives
\(\partial_4\operatorname{pr}_4=\Sigma(W\to F_3)d\).
Restricting to the summand \(U_4\) gives the displayed formula.

We next express \(d\) through the height-two fracture map.  Choose
coherent inverse data for \(L_{K(2)}w\), and form
\[
 \tau_2:L_1X\longrightarrow L_1L_{K(2)}X
       \xrightarrow{(L_1L_{K(2)}w)^{-1}}L_1L_{K(2)}W.
\]
It is a map under \(L_1W\), so it induces
\[
 \bar\tau_2:Y\simeq\cofib(L_1w)
       \longrightarrow\cofib(L_1W\to L_1L_{K(2)}W).
\]
Let \(j_W:\fib(W\to L_{K(2)}W)\to W\) be the fiber inclusion.
Then
\begin{equation}\label{eq:first-roof}
\begin{split}
 d:\ U_{\mathrm{mid}}\xrightarrow hY\xrightarrow{\bar\tau_2}
 &\cofib(L_1W\longrightarrow L_1L_{K(2)}W)\\
 &\simeq\Sigma\fib(W\longrightarrow L_{K(2)}W)
   \xrightarrow{-\Sigma j_W}\Sigma W.
\end{split}
\end{equation}
Indeed, \eqref{eq:height-two-fracture} identifies the fiber of
\(W\to L_{K(2)}W\) with the fiber of
\(L_1W\to L_1L_{K(2)}W\); it is \(E(1)\)-local.
\Cref{lem:boundary}, with \(L=L_{K(2)}\) and \(M=L_1\),
then says that the composite from \(Y\) in
\eqref{eq:first-roof} is exactly \(\partial_w\).  This proves
the formula and fixes its minus sign.

\emph{The rational quotient.}
Let \(a:F_4\to X\) be the pullback projection and set
\(Z=\cofib(a)\).  Taking the cofiber of the pullback square defining
\(F_4\) gives a natural equivalence
\begin{equation}\label{eq:quotient-identification}
 Z\simeq\cofib(h:U_{\mathrm{mid}}\longrightarrow Y).
\end{equation}
This cofiber is \(E(1)\)-local and \(K(1)\)-acyclic by
\eqref{eq:middle-map}, hence rational.  Thus \(a\) is an
equivalence after localization at \(K(1)\) or \(K(2)\), and
therefore after \(L_J\).  Rationalizing
\eqref{eq:quotient-identification} identifies
\(Z\simeq\cofib(r)\).

Write \(z:X\to Z\) for the quotient.  Since \(Z\) is rational,
\eqref{eq:rational-adjunction} gives its factorization
\(\widetilde z:L_0X\to Z\).  Under the preceding cofiber
identification, the map in \textup{(iii)} is
\begin{equation}\label{eq:terminal-quotient}
 t=\widetilde z\,b:V\xrightarrow{\simeq}Z.
\end{equation}
This records the factorization used to compare the last attachments
with fracture.

\emph{The last four stages.}
Let \(V_{\leq j}\) be the first \(j\) summands of \(V\), for
\(0\leq j\leq4\), and use the restriction of \(t\) to define
\begin{equation}\label{eq:last-stages}
 F_{4+j}=X\mathbin{\times}_ZV_{\leq j}.
\end{equation}
At \(j=0\), the cofiber sequence of \(a\) identifies this
definition with the previous \(F_4\), over \(X\).
The cofiber of \(V_{\leq j-1}\to V_{\leq j}\) is \(V_{4+j}\).
Exact pullback gives the same cofiber for
\(F_{4+j-1}\to F_{4+j}\).

Let \(\partial_a:Z\to\Sigma F_4\) be the boundary of \(a\),
put \(g=\partial_at\), and write \(g_i=g|_{V_i}\).
The quotient-of-a-pullback argument used above gives every remaining
connecting map:
\begin{equation}\label{eq:last-boundaries}
 \partial_i=\Sigma(F_4\to F_{i-1})g_i:
 V_i\longrightarrow\Sigma F_{i-1},\qquad 5\leq i\leq8.
\end{equation}

The map \(g\) is described by the positive-height fracture square.
Choose coherent inverse data for \(L_Ja\), and let
\(j_a:\fib(F_4\to L_JF_4)\to F_4\) be the fiber inclusion.
The formula is
\begin{equation}\label{eq:second-roof}
\begin{aligned}[b]
 g:\ V\xrightarrow bL_0X\longrightarrow L_0L_JX
 &\xrightarrow{(L_0L_Ja)^{-1}}L_0L_JF_4\\
 &\longrightarrow\cofib(L_0F_4\longrightarrow L_0L_JF_4)\\
 &\simeq\Sigma\fib(F_4\longrightarrow L_JF_4)
   \xrightarrow{-\Sigma j_a}\Sigma F_4.
\end{aligned}
\end{equation}
We verify in particular why this roof begins with \(b\), rather
than only with the quotient equivalence \(t\).
Put
\[
 C_a=\cofib(F_4\to L_JF_4),\qquad
 \tau_J=(L_Ja)^{-1}\circ(X\to L_JX),
\]
and let \(q_a:L_JF_4\to C_a\) be the quotient.  The chosen inverse
homotopy makes \(\tau_Ja\) the localization unit of \(F_4\).
\Cref{lem:boundary} gives \(\bar\tau_J:Z\to C_a\), with a
specified homotopy
\[
 \bar\tau_Jz\simeq q_a\tau_J.
\]
Both \(Z\) and \(C_a\) are rational: for \(C_a\) this follows
from \eqref{eq:positive-fracture}.  If
\(\rho:C_a\xrightarrow{\simeq}L_0C_a\) is its localization map,
the rational adjunction transports the last homotopy to
\[
 \rho\bar\tau_J\widetilde z\simeq (L_0q_a)(L_0\tau_J).
\]
Using \(t=\widetilde z b\), this becomes
\[
 \rho\bar\tau_Jt\simeq
 (L_0q_a)(L_0L_Ja)^{-1}\circ
 (L_0X\to L_0L_JX)\circ b.
\]
Exactness identifies \(L_0C_a\) with the cofiber in
\eqref{eq:second-roof}.  The boundary formula
\(\partial_a=(-\Sigma j_a)\bar\tau_J\) now identifies that
entire roof with \(\partial_at=g\), including the second minus sign.

Finally, \eqref{eq:terminal-quotient} makes the projection
\[
 e:F_8=X\times_ZV\longrightarrow X
\]
an equivalence.  All stages use finite limits or colimits of
\(E(2)\)-local \(S\)-modules, so they remain in that category.
Their maps to \(X\) extend \(w\); hence
\(F_2\to F_8\xrightarrow eX\) is \(w\).
\end{proof}

\begin{theorem}[Attachment identification for the ordered fracture filtration]
\label{thm:attachment-identification}
Fix the data \(w,\kappa,r,b\) and the compatibility homotopy of
\Cref{prop:fracture}.  Choose coherent inverse data for
\(L_{K(2)}w\), and carry out that proposition's construction.
Use exactly its fracture map \(h\), pullback stages \(F_i\),
projection \(a:F_4\to X\), quotient \(Z=\cofib(a)\),
equivalence \(t=\widetilde z b:V\xrightarrow{\simeq}Z\), and
terminal projection \(e:F_8\to X\), choosing the coherent inverse
data for \(L_Ja\) at the indicated step.  Under the resulting
ordered cofiber identifications
\[
 W_1,\ W_2,\ U_3,\ U_4,\ V_5,\ V_6,\ V_7,\ V_8,
\]
the connecting maps of this same filtration are
\[
\begin{gathered}
 \partial_1=0,\qquad \partial_2=0,\qquad
 \partial_3=d_3,\qquad
 \partial_4=\Sigma(F_2\to F_3)d_4,\\
 \partial_i=\Sigma(F_4\to F_{i-1})g_i
 \qquad(5\leq i\leq8),
\end{gathered}
\]
where \(d=\partial_wh\), \(g=\partial_at\),
\(d_i=d|_{U_i}\), and \(g_i=g|_{V_i}\).  The maps \(d\) and \(g\)
are the signed roofs
\[
\begin{aligned}
 d:\ U_{\mathrm{mid}}\xrightarrow hY\xrightarrow{\bar\tau_2}
 &\cofib(L_1W\longrightarrow L_1L_{K(2)}W)\\
 &\simeq\Sigma\fib(W\longrightarrow L_{K(2)}W)
   \xrightarrow{-\Sigma j_W}\Sigma W,
\end{aligned}
\]
and
\[
\begin{aligned}
 g:\ V\xrightarrow bL_0X\longrightarrow L_0L_JX
 &\xrightarrow{(L_0L_Ja)^{-1}}L_0L_JF_4\\
 &\longrightarrow\cofib(L_0F_4\longrightarrow L_0L_JF_4)\\
 &\simeq\Sigma\fib(F_4\longrightarrow L_JF_4)
   \xrightarrow{-\Sigma j_a}\Sigma F_4.
\end{aligned}
\]
Here \(\bar\tau_2,j_W,j_a\) are the maps defined from the same
construction and inverse data above.
\end{theorem}

\begin{proof}
The proof of \Cref{prop:fracture} established
\eqref{eq:first-boundaries} and \eqref{eq:last-boundaries} under
these ordered cofiber identifications.  It identified their
components by \eqref{eq:first-roof} and \eqref{eq:second-roof}
using precisely the stated inverse data.  The collected formulas
therefore apply to the same \(h,F_i,a,Z,t,e\).
\end{proof}

\begin{remark}
The choices of compatibility homotopy and coherent inverses are
part of this realization.  \Cref{thm:attachment-identification}
does not claim that the
filtration is unique or that its attachments are canonical without
those choices.  Together with \Cref{prop:fracture}, it identifies the
original map \(w\), the terminal projection, and the connecting maps
associated with the stated data.
In the height-three application, only the first component of \(w\)
is asserted to be the canonical unit.
\end{remark}
